\section{Joining matrix cones and localizing traces}
\label{sec:localization}

Our goal is to retain the tracial mass from the previous section while
obtaining norm-central maps. The trace-kernel ideal and relative
commutant lifting methods have predecessors in Kirchberg--R\o rdam
\cite[Section~4]{KirchbergRordam2014}. Here the relative ideal is
specified by hereditary rank tests, including the traces which vanish
on the controlling positive element.

Throughout this section, traces on $A$ are the lower semicontinuous
extended traces of the preliminaries.  In particular, they need not be
finite on a dense ideal.  If $D\subseteq A$ is hereditary, write
$T_{\leq1}(D)$ for the bounded positive tracial functionals on $D$ of
norm at most one, and put
\[
 \|z\|_{2,T_{\leq1}(D)}
   =\sup_{\sigma\in T_{\leq1}(D)}\sigma(z^*z)^{1/2}.
\]
The supremum is zero if there are no nonzero such functionals.  Expressions
involving $1-h$ are evaluated in a unitization; a sandwich by an element
of $D$ still belongs to $D$.

\begin{lemma}[Joining two matrix cones]\label{lem:join}
Let $B$ be a $C^*$-algebra, and let $\phi_1,\phi_2:M_k\to B$ be
\cpc order-zero maps with commuting ranges.  Set $h_i=\phi_i(1)$.
There is a \cpc order-zero map $\phi:M_k\to
C^*(\phi_1(M_k),\phi_2(M_k))$ such that
\begin{equation}\label{loc:join-weight}
 \phi(1)=1-(1-h_1)(1-h_2)=h_1+h_2-h_1h_2.
\end{equation}
Its weight dominates both $h_1$ and $h_2$.  If both original maps commute
with a subset of an ambient algebra, the joined map also commutes with
that subset.
\end{lemma}

\begin{proof}
Let $\mathcal C_k=C_0((0,1],M_k)$. Its unitization is the algebra of continuous
$M_k$-valued functions on $[0,1]$ whose value at zero is scalar.  Realize
$\widetilde{\mathcal C}_k\otimes\widetilde{\mathcal C}_k$ as the algebra of continuous
$M_k\otimes M_k$-valued functions $F$ on the square satisfying
\[
 F(x,0)\in M_k\otimes1,
 \qquad F(0,y)\in1\otimes M_k.
\]
These conditions also make $F(0,0)$ scalar.  Here is a direct verification
that these conditions characterize the tensor product.  The tensor
product contains every scalar continuous function on the square, and
its evaluation at each point is the entire indicated fiber: use scalar
functions vanishing at zero to realize matrix entries on either
nonzero coordinate.  Given a continuous function satisfying the
conditions, choose, at each point, an element of the tensor product
having the same value there.  Continuity gives a neighborhood on which
the error is less than a prescribed tolerance.  A finite scalar
partition of unity subordinate to these neighborhoods gives an
approximation in the tensor product.  Closedness proves the assertion.

Let $S$ be the unitary that interchanges the two factors of
$\C^k\otimes\C^k$, and put $Q=(1-S)/2$.  Since $S=S^*=S^{-1}$,
$Q$ is a projection, and
\[
 U(t)=\exp(i\pi tQ),\qquad 0\le t\le1,
\]
is a unitary path from $1$ to $S$.  For $z\in M_k$ define
\[
 F_z(x,y)=
 \begin{cases}
 (x+y-xy)\,U\bigl(y/(x+y)\bigr)(z\otimes1)
                 U\bigl(y/(x+y)\bigr)^*, &(x,y)\ne(0,0),\\
 0,&(x,y)=(0,0).
 \end{cases}
\]
At $y=0$ this is $xz\otimes1$, and at $x=0$ it is $y1\otimes z$.
At the origin continuity follows from
$\|F_z(x,y)\|\le(x+y)\|z\|$; no limit of the unitary itself is needed.
Thus $F_z$ belongs to the tensor product.  Pointwise it is a positive
scalar contraction times a $*$-homomorphism, so $z\mapsto F_z$ is
completely positive, contractive, and order zero.  Its weight is
$(x+y-xy)1$.

The cone homomorphisms associated to $\phi_1$ and $\phi_2$ extend to
commuting unital homomorphisms from their unitizations.  Nuclearity of
the cones lets their product be evaluated on the spatial tensor
product.  The functions $F_z$ vanish at $(0,0)$.  The ideal of functions
vanishing there is generated by $\mathcal C_k\otimes1$ and $1\otimes\mathcal C_k$:
quotienting by those two ideals leaves only the scalar evaluation at
the origin.  Its image consequently lies in
$C^*(\phi_1(M_k),\phi_2(M_k))$, even when $B$ is nonunital.  Evaluation
of $F_z$ gives the asserted map and its weight.  Finally
$\phi(1)-h_1=(1-h_1)h_2\ge0$, and similarly for $h_2$.
The last assertion follows from containment of the range in the
generated algebra.
\end{proof}

\begin{lemma}[Filling bounded traces]\label{loc:bounded-fill}
Let $D\subseteq A$ be hereditary, $k\ge2$, $F\subset D$ finite, and
$\mathcal P\subset D$ a finite set of positive contractions.  For every
$\rho>0$ there is a \cpc map $\mu:M_k\to D$ such that, for all matrix
units $e_{ij},e_{rs}$ and all $x\in F$,
\begin{align}
 \|[\mu(e_{ij}),x]\|_{2,T_{\leq1}(D)}&<\rho,\label{loc:fill-comm}\\
 \|\mu(e_{ij})\mu(e_{rs})-\mu(1)\mu(e_{ij}e_{rs})
       \|_{2,T_{\leq1}(D)}&<\rho,\label{loc:fill-oz}\\
 \sigma\bigl(a^{1/2}(1-\mu(1))a^{1/2}\bigr)&<\rho
       \quad(a\in\mathcal P,\ \sigma\in T_{\leq1}(D)).
       \label{loc:fill-trace}
\end{align}
\end{lemma}

\begin{proof}
The sequences that vanish in the uniform $2$-seminorm define a closed
two-sided ideal $J_2$ of $D_\omega$.  Indeed the uniform $2$-seminorm is
at most the norm, and
$\|uz\|_2,\|zu\|_2\le\|u\|\|z\|_2$ for every bounded trace; the
triangle inequality and the same inequalities establish ideality and
closedness.  Write $B=D_\omega/J_2$.  Every choice
$\sigma_n\in T_{\leq1}(D)$ gives a bounded trace
\[
 \sigma_\omega((z_n))=\lim_{n\to\omega}\sigma_n(z_n)
\]
on $D_\omega$, which factors through $B$ by Cauchy--Schwarz.  These
choices are allowed to depend on $n$.

Fix $0<\delta<1$.  Proposition~\ref{prop:bounded-trace}, applied at
coordinate $n$ with its auxiliary error tending to zero, has the
following consequence.  Given finitely many positive contractions in
$B$ and finitely many elements with which to commute, there is a
\cpc order-zero map $\psi:M_k\to B$ commuting with those elements and
satisfying
\begin{equation}\label{loc:fraction}
 \sigma_\omega(c\psi(1))
 \ge\frac{\delta}{3}\sigma_\omega(c)-C_k\delta^{3/2}
\end{equation}
for all the chosen targets and all the traces just described.
To see the quantifiers, choose positive contractive lifts of the
finitely many targets and bounded lifts of the tests, and apply that
proposition separately at every coordinate, including real and
imaginary parts of the tests, rescaled to contractions, with the
requested errors tightened by their fixed norms.  The coordinate errors tend to zero
uniformly over $T_{\leq1}(D)$.  The coordinate maps are order zero,
and their commutator sequences belong to $J_2$.  This proves the
assertion.  The lifts exist by completely positive lifting for the
finite-dimensional domain when a map, rather than a single element,
is being lifted.

Start with joined weight $H_0=0$.  Having constructed the preceding
joined map, let $r=1-H_j$ and prescribe the targets
$r^{1/2}ar^{1/2}$ for $a\in\mathcal P$.  Prescribe commutation with
the preceding map and with the fixed tests.  Apply
\eqref{loc:fraction} and join the new map to the preceding one using
Lemma~\ref{lem:join}.  The new weight commutes with $r$, and hence, for
any trace $\sigma_\omega$, the defect
\[
 D_j(a)=\sigma_\omega\bigl(a^{1/2}(1-H_j)a^{1/2}\bigr)
\]
satisfies
\[
 D_{j+1}(a)\le(1-\delta/3)D_j(a)+C_k\delta^{3/2}.
\]
This follows by subtracting
$\sigma_\omega(r^{1/2}ar^{1/2}\psi(1))$ from $D_j(a)$;
traciality and commutation of $\psi(1)$ with $r$ justify the subtraction
without any commutation assumption on $a$.  Since $D_0(a)\le1$, induction
gives
\begin{equation}\label{loc:geometric}
 D_N(a)\le(1-\delta/3)^N+3C_k\sqrt\delta.
\end{equation}
Choose $\delta$ so that the second term is less than $\rho/4$, and
then a finite $N$ so that the first term is less than $\rho/4$.

Lift the resulting \cpc map from $B$ to coordinate \cpc maps into $D$.
Its exact order-zero and commutation identities in $B$ give
\eqref{loc:fill-comm}--\eqref{loc:fill-oz} on an $\omega$-large set.
Its trace estimates also give \eqref{loc:fill-trace} on an
$\omega$-large set.  Otherwise choose, at every failing coordinate, a
violating tracial contraction and one of the finitely many targets;
pass to an $\omega$-large set where the target is fixed.  Their
ultralimit would contradict the strict bound $\rho/2$ from
\eqref{loc:geometric}.  Intersect these finitely many good sets and
choose a coordinate.  This proves all the required uniform estimates.
\end{proof}

We now seek norm-central order-zero maps whose trace deficits against
a positive contraction $a$ are arbitrarily small relative to the rank
of a fixed $b\in A_+$. Here $a\in\Ideal_A(b)$ and
$d_\tau(a)\le C d_\tau(b)$ for every $\tau$ and one finite $C>0$; the precise
statement is Proposition~\ref{prop:localized-trace} below. For positive
contractive representatives $(h_n)$ of the weight, the target is
\[
 \tau\bigl(a^{1/2}(1-h_n)a^{1/2}\bigr)\le\eta d_\tau(b)
\]
on a common $\omega$-large set for each $\eta>0$, uniformly over traces
with $d_\tau(b)<\infty$. To remove the norm defects without losing this
estimate, we measure an error by the rank of each positive cut of its modulus.
Taking cuts makes this smallness condition stable under norm-null
changes of representatives; traces with $d_\tau(b)=0$ will require
exact vanishing.

\begin{definition}\label{loc:relative-ideal}
For $b\in A_+$ define $J_b\subseteq A_\omega$ as follows.  An element
$z=(z_n)$ belongs to $J_b$ if, for every $s>0$ and $\eta>0$, there is
$E\in\omega$ such that
\begin{equation}\label{loc:J-predicate}
 d_\tau((|z_n|-s)_+)\le\eta d_\tau(b)
 \quad\text{for every }n\in E\text{ and every }\tau
          \text{ with }d_\tau(b)<\infty.
\end{equation}
For $d_\tau(b)=0$ this means exact zero on the left.  The set $E$ is
common to the entire indicated trace family.  It is enough to test
positive rational $s,\eta$.
\end{definition}

\begin{lemma}[The relative trace ideal]\label{loc:sigma-ideal}
Definition~\ref{loc:relative-ideal} is independent of the bounded
representative.  It defines a closed two-sided ideal of $A_\omega$.
Moreover, for every separable $C^*$-subalgebra $C_0\subseteq A_\omega$
there is a positive contraction $e\in J_b\cap C_0'$ such that
$ey=y$ for every $y\in J_b\cap C_0$.
\end{lemma}

\begin{proof}
We record the rank facts used in verifying ideality.  For $u,v,x\in A$,
\begin{equation}\label{loc:rank-calculus}
 d_\tau(|u^*|)=d_\tau(|u|),\quad
 d_\tau(|xu|),d_\tau(|ux|)\le d_\tau(|u|),\quad
 d_\tau(|u+v|)\le d_\tau(|u|)+d_\tau(|v|).
\end{equation}
The first follows from Cuntz equivalence of the two squared moduli.
For the second use $u^*x^*xu\le\|x\|^2u^*u$, and then the first
identity for right multiplication.  For the last use
$(u+v)^*(u+v)\le2(u^*u+v^*v)$ and the Cuntz inequality for a sum of
positive elements.

For $t>0$ the polar cut
\[
 u^{[t]}=u\,q_t(|u|),\qquad
 q_t(\lambda)=
 \begin{cases}(\lambda-t)_+/\lambda,&\lambda>0,\\0,&\lambda=0,
 \end{cases}
\]
belongs to $A$, has modulus $(|u|-t)_+$, and satisfies
$\|u-u^{[t]}\|\le t$.  Consequently a fixed positive cut of a sum or
product can be Cuntz compared to the modulus of a sum or product of
sufficiently small polar cuts.  In detail, choose $t$ so small that
the norm difference of the two expressions makes the difference of
their moduli less than the desired cut.  This is possible uniformly
on the bounded sets at issue, by uniform continuity of
$z\mapsto|z|$.  Norm comparison and \eqref{loc:rank-calculus} then bound
that cut rank by the sum of the ranks of the input polar cuts, or by
one such rank for multiplication and adjoints.  Apply
\eqref{loc:J-predicate} with error $\eta/2$ for each of the two
summands, and with error $\eta$ for a product or adjoint.  Finite
intersection of the good sets proves closure under addition,
adjoints, and left and right multiplication by bounded sequences.
Scalar multiplication follows in the same way.

Here are the perturbation details, including the zero-rank case.
If $(z_n)$ and $(z'_n)$ represent the same element, then, for a fixed
$s>0$, on an $\omega$-large set
$\||z'_n|-|z_n|\|<s/2$.  Therefore
\begin{equation}\label{loc:representative-cut}
 [ (|z'_n|-s)_+ ]\le[ (|z_n|-s/2)_+ ],
\end{equation}
by comparing $|z'_n|$ with $(|z_n|-s/2)_+$ in norm, at error strictly
less than $s$.  This proves representative independence.  If
$z^{(l)}\in J_b$ converges in the norm of $A_\omega$ to $z$, fix a cut
$s>0$ and choose $l$ so that bounded lifts of $z^{(l)}$ and $z$ are
close enough on an $\omega$-large set to give
\eqref{loc:representative-cut} with $z^{(l)}$ on the right.  Apply its
predicate at cut $s/2$ and arbitrary error $\eta$.  This proves norm
closedness.  When $d_\tau(b)=0$, every rank comparison in these
arguments has an exactly zero right side on the common good set.
No assertion about the uncut rank of a norm-small perturbation is
being used.

Put $I_0=J_b\cap C_0$.  Choose dense sequences $(c_j)$ in the unit ball
of $C_0$ and $(y_j)$ in the unit ball of $I_0$.  A quasicentral
approximate unit for $I_0$ in $C_0$ supplies positive contractions
$e^{(l)}\in I_0$ such that, for $j\le l$,
\[
 \|[e^{(l)},c_j]\|<1/l,\qquad
 \|e^{(l)}y_j-y_j\|<1/l.
\]
Choose positive contractive lifts $e^{(l)}_n$ and bounded lifts of the
$c_j,y_j$.  For each $l$, there is an $\omega$-large set $G_l$ on
which these finitely many norm estimates hold with $2/l$ in place of
$1/l$ and, simultaneously, for $1\le j\le l$,
\begin{equation}\label{loc:sigma-stage}
 d_\tau((e^{(l)}_n-1/j)_+)\le l^{-1}d_\tau(b)
       \quad\text{for all }\tau\text{ with }d_\tau(b)<\infty.
\end{equation}
Each displayed inequality is one uniform coordinate predicate, so
only finitely many predicates are intersected at this stage.  In
particular, this construction does not enumerate traces.

Replace $G_l$ by the decreasing sets
$E_l=\{n\ge l\}\cap\bigcap_{r\le l}G_r$.
For $n\in E_1$ let $l(n)$ be the largest $l\le n$ with $n\in E_l$,
and set $e_n=e^{(l(n))}_n$; put $e_n=0$ elsewhere.  Then
$l(n)\to_\omega\infty$.  The norm tests imply that $e=(e_n)$
commutes with $C_0$ and absorbs $I_0$.  Given $s,\eta>0$, choose
$j$ with $1/j<s$ and then require $l(n)\ge j$ and $1/l(n)\le\eta$.
Equation~\eqref{loc:sigma-stage} implies
\eqref{loc:J-predicate}.  For zero-rank traces it continues to assert
exact zero at each selected coordinate.  Thus $e\in J_b$, completing
the proof.
\end{proof}

\begin{proposition}[Localized norm-central trace filling]
\label{prop:localized-trace}
Let $k\ge2$, $b\in A_+$, and let $a\in\Ideal_A(b)$ be a positive
contraction.  Suppose that, for some finite $C>0$,
\begin{equation}\label{loc:relative-rank}
 d_\tau(a)\le C d_\tau(b)\qquad\text{for every }\tau.
\end{equation}
There is a \cpc order-zero map $\phi:M_k\to A_\omega\cap A'$ with
weight $h$ such that every positive contractive representative
$(h_n)$ of $h$ has the following property: for every $\eta>0$, on a
common $\omega$-large set of coordinates,
\begin{equation}\label{loc:TL}
 \tau\bigl(a^{1/2}(1-h_n)a^{1/2}\bigr)
       \le\eta d_\tau(b)
 \qquad(d_\tau(b)<\infty).
\end{equation}
In particular the left side is exactly zero when $d_\tau(b)=0$.

The finite-coordinate form permits arbitrary relative trace accuracies
together with prescribed finite norm tests. Given a
finite set $F\subset A$, finitely many pairs $(a_l,b_l)$ satisfying
\eqref{loc:relative-rank} with finite constants $C_l$, positive
accuracies $\eta_l$, and $\rho>0$, there is a \cpc map
$\theta:M_k\to A$ whose matrix-unit order-zero defects and
commutators with $F$ have norm less than $\rho$, and whose weight
satisfies
\begin{equation}\label{loc:TL-coordinate}
 \tau\bigl(a_l^{1/2}(1-\theta(1))a_l^{1/2}\bigr)
     \le\eta_l d_\tau(b_l)
 \qquad(d_\tau(b_l)<\infty)
\end{equation}
for every $l$, simultaneously over all the indicated traces.
\end{proposition}

\begin{proof}
The relative rank bound will control both the approximate-unit tail
in the hereditary construction and the trace mass lost when we remove
the ideal $J_b$.
Write $I=\Ideal_A(b)$ and choose a sequential positive contractive
quasicentral approximate unit $v_i$ for $I$ in $A$.  Fix $t_i>0$
tending to zero, and set
\[
 p_i=(v_i-2t_i)_+,\qquad q_i=\chi_i(v_i),
\]
where $\chi_i$ is zero on $[0,t_i]$, increases linearly to one on
$[t_i,2t_i]$, and is one thereafter.  We may assume $2t_i<1$.
Then $p_iq_i=p_i$, $\|p_i-v_i\|\le2t_i$, and $(p_i)$ is still a
quasicentral approximate unit.  Moreover
$[q_i]\le[(v_i-t_i/2)_+]$.  Compact containment and $v_i\in I$ give
an integer $L_i\ge1$ with
\begin{equation}\label{loc:hereditary-bound}
 [q_i]\le L_i[b].
\end{equation}
No uniform bound on these integers is required.

Let $D_i=\her_A(q_i)$.  If $d_\tau(b)<\infty$, the restriction of
$\tau$ to $D_i$ is a bounded positive tracial functional with
\begin{equation}\label{loc:trace-norm}
 \|\tau|_{D_i}\|=d_\tau(q_i)\le L_i d_\tau(b).
\end{equation}
Indeed every positive contraction of $D_i$ is supported by $p_{q_i}$,
so its trace is at most $d_\tau(q_i)$; an approximate unit obtained
from $q_i$ shows equality.  Equivalently, use $q_i^{1/m}$ and lower
semicontinuity.  Finiteness of the bound turns the additive tracial
map on positives into a bounded positive linear functional on $D_i$.
For $d_\tau(b)=0$ that functional is zero.  For a positive finite
$d_\tau(b)$, dividing it by $L_i d_\tau(b)$ gives an element of
$T_{\leq1}(D_i)$.

Choose a dense sequence $(x_j)$ in the unit ball of $A$, and apply
Lemma~\ref{loc:bounded-fill} inside $D_i$, with tests
\[
 p_i,\quad q_i x_j q_i\ (j\le i)
 \quad\text{and their real and imaginary parts},
\]
and positive target $p_iap_i$.  Choose its accuracy
$\rho_i>0$ after $L_i$ has been fixed, so that
\begin{equation}\label{loc:accuracy}
 L_i\rho_i\longrightarrow0,
 \qquad L_i\rho_i^2\longrightarrow0;
\end{equation}
for example $\rho_i=1/(i(1+L_i))$ works.  Denote the resulting maps
by $\mu_i:M_k\to D_i$, and put
$\theta_i(z)=p_i\mu_i(z)p_i$.  These are \cpc maps into $A$.

Their trace defects have the required relative bound already before
correction.  Put
$\lambda_i=\|a^{1/2}(1-p_i^2)a^{1/2}\|$, so $\lambda_i\to0$.
The positive element defining $\lambda_i$ is supported by $p_a$;
using \eqref{loc:relative-rank}, its trace is at most
$C\lambda_i d_\tau(b)$.  The other part of the defect is, by
traciality,
\[
 \tau\bigl((1-\mu_i(1))^{1/2}p_iap_i
                   (1-\mu_i(1))^{1/2}\bigr)
 \le L_i\rho_i d_\tau(b).
\]
Thus
\begin{equation}\label{loc:uncorrected-trace}
 \tau\bigl(a^{1/2}(1-\theta_i(1))a^{1/2}\bigr)
 \le(C\lambda_i+L_i\rho_i)d_\tau(b)
 \qquad(d_\tau(b)<\infty).
\end{equation}
At zero rank both contributions are exactly zero.

Let $\mu,\theta$ be the \cpc maps represented by the coordinate
maps.  Their relevant relations first hold modulo $J_b$.  In fact,
for each tested $2$-norm defect $z_i$, normalization in
\eqref{loc:trace-norm} and functional calculus give
\[
 d_\tau((|z_i|-s)_+)
   \le s^{-2}\tau(z_i^*z_i)
   \le s^{-2}L_i\rho_i^2d_\tau(b).
\]
The estimate holds for every fixed $s>0$; at zero rank it holds because
all the coordinates at issue belong to $I$ and the restricted trace
is zero.  Equation~\eqref{loc:accuracy} proves that these defect
sequences belong to $J_b$.  Consequently, modulo $J_b$, the map
$\mu$ is order zero and commutes with $p=(p_i)$ and each
$(q_i x_jq_i)$.  Also
$[p_i,q_ix_jq_i]=q_i[p_i,x_j]q_i\to0$ in norm.  It follows that
$\theta=p\mu(\cdot)p$ is order zero modulo $J_b$ and commutes there
with $(q_ix_jq_i)$.  It commutes there with the constant $x_j$ as
well: using $q_ip_i=p_i$, the difference between the two commutators
has norm at most $2\|[p_i,x_j]\|$.  For example
$p_i\mu_i(z)p_i x_j(1-q_i)
=p_i\mu_i(z)[p_i,x_j](1-q_i)$ for $\|z\|\le1$; the other term is
its analogous left-sided expression.  Density gives the assertion
for all of $A$.

Apply Lemma~\ref{loc:sigma-ideal} to
$C_0=C^*(A,\theta(M_k))$, and let $e$ be its absorbing positive
contraction.  Define
\[
 \phi(z)=(1-e)\theta(z).
\]
Since $e$ commutes with the range, this is c.p.c.  Every order-zero
polynomial defect of $\theta$ belongs to $J_b\cap C_0$ and is
annihilated by $1-e$.  Multiplication by $(1-e)^2$ therefore proves
the exact order-zero identities for $\phi$.  Likewise
$[\phi(z),x]=(1-e)[\theta(z),x]=0$ for $x\in A$.  Thus
$\phi:M_k\to A_\omega\cap A'$ has the required norm relations.

To check the trace loss, the removed weight
$j=e\theta(1)$ is a positive contraction in $J_b$.  Choose positive
contractive representatives $j_i$ and $h_i$ of $j$ and $h=\phi(1)$.
Then $\theta_i(1)-h_i-j_i\to_\omega0$ in norm.  For $s>0$,
functional calculus gives $j_i\le s1+(j_i-s)_+$, whence
\begin{equation}\label{loc:correction-cost}
 \tau(a^{1/2}j_i a^{1/2})
 \le s d_\tau(a)+d_\tau((j_i-s)_+).
\end{equation}
Here the second term bounds the trace of a compression of the
positive contraction $(j_i-s)_+$, and its rank is unchanged by a
positive rescaling if necessary.  Given a desired relative error,
first choose $s$ so that $sC$ is small, then use
\eqref{loc:J-predicate} at that fixed cut with a sufficiently small
error.  Norm-null differences of the lifts contribute at most their
norm times $d_\tau(a)\le C d_\tau(b)$.  Combining these observations
with \eqref{loc:uncorrected-trace} proves \eqref{loc:TL}, uniformly
for positive finite $d_\tau(b)$.  If $d_\tau(b)=0$, then
$d_\tau(a)=0$ and every positive contraction $h_i$ has exactly zero
defect against $a$.  This also proves the assertion for every choice
of positive contractive representative of $h$.

For completeness, the passage to several pairs permits arbitrary
finite coordinate tests.  Construct the map for the first pair.
For the next pair use the preceding construction and reindex its
coordinate lifts as follows.  At coordinate $n$ the matrix-unit
values of all already fixed maps are a finite list of ordinary
elements of $A$.  Norm centrality of the new map allows an index at
which its commutators with that list and with the first $n$ fixed
dense tests are less than $1/n$.  Intersect this set of indices with
the sets imposing its order-zero errors and its uniform trace
predicate for its fixed pair, to any desired accuracy tending to
zero.  A finite intersection is $\omega$-large and hence nonempty.
The selected sequence still has \eqref{loc:TL} and now commutes with
the earlier maps.  Repeat for the finite list of pairs and join.

The joined weight $H$ dominates every original weight $h^{(l)}$.
This domination also preserves the uniform trace estimates at
coordinates.  Indeed choose positive lifts $D^{(l)}_n$ of
$H-h^{(l)}$; for any positive contractive lifts $H_n,h^{(l)}_n$,
\[
 H_n\ge h^{(l)}_n-\delta_{l,n}1,
 \qquad \delta_{l,n}\to_\omega0.
\]
Their trace-defect difference against $a_l$ is therefore at most
$C_l\delta_{l,n}d_\tau(b_l)$ when that rank is finite; at zero rank
both defects vanish.  Lift the joined \cpc map and select one
coordinate in the finite intersection imposing all desired norm
relations, all prescribed commutators, and all relative trace
bounds with strict slack.  This gives
\eqref{loc:TL-coordinate}.  The finite set $F$, and each $\eta_l$,
are arbitrary at this final step.
\end{proof}
